% ============================================================
%  Chapter 2 — Literature Review
% ============================================================
\chapter{Literature Review}
\label{ch:literature}

\section{Overview}
\label{sec:lit_overview}

The study of where groundwater accumulates, and how to map it from the surface,
sits at the intersection of hydrogeology, remote sensing, and---more recently---
machine learning. This chapter traces that intersection. It begins with the
geo-environmental factors that the literature has settled on as the principal
controls on groundwater occurrence, since these are the raw material of every
potential map. It then follows the methodological arc that GWPM research has
taken over the last two decades: from knowledge-driven weighted overlay and the
analytic hierarchy process, through bivariate and multivariate statistical
techniques, to the tree-based machine-learning models that now dominate spatial
prediction. A separate section examines the broader uptake of machine learning
in hydrology, which supplies both the methods and the justification for applying
them here. The chapter closes by drawing together what has, and has not, been
done in the \studyarea{} and its surrounding region, which is where the gap this
thesis addresses comes into focus.

\section{Geo-environmental Controls on Groundwater}
\label{sec:lit_factors}

There is broad agreement in the literature about which surface and near-surface
variables govern groundwater occurrence, even if the exact list varies from one
study to the next. Geology and soil set the stage: lithology determines the
storage and transmission properties of the subsurface, while soil texture
controls how readily rainfall infiltrates rather than running off. Terrain
matters through slope and the drainage network---gentle slopes and low-lying,
convergent areas favour accumulation, steep slopes favour runoff. Climate enters
chiefly through rainfall, the ultimate source of recharge. Vegetation and
moisture, captured through spectral indices, and surface temperature, a proxy
for evapotranspiration, round out the picture.

Arulbalaji et al. \cite{arulbalaji2019gisahp} offer a representative example of
how comprehensive these factor stacks can become, integrating twelve thematic
layers---geology, geomorphology, land use, lineament density, drainage, rainfall,
soil, slope, roughness, topographic wetness index (TWI), topographic position,
and curvature---in their assessment of the Western Ghats. Most studies work with
a smaller but overlapping subset. The topographic wetness index, introduced by
Beven and Kirkby \cite{beven1979topmodel} as part of the TOPMODEL framework,
recurs across the GWPM literature because it compactly expresses the tendency of
a location to accumulate water based on its upslope contributing area and local
slope. Morphometric analyses of drainage basins reinforce the same theme, linking
network properties such as drainage density and stream order to infiltration and
groundwater prospects \cite{mahala2019morphometric,kumar2018morphometric}.

The practical upshot for this thesis is that the nine features adopted here---
elevation, slope, aspect, TWI, NDVI, NDWI, rainfall, soil texture, and land
surface temperature---are not an arbitrary selection but a subset that recurs
throughout the GWPM literature, chosen for being derivable consistently from
freely available remote-sensing archives across the whole plateau.

\section{Knowledge-Driven Methods: Weighted Overlay and AHP}
\label{sec:lit_ahp}

For most of the field's history, the standard way to turn a stack of factor
layers into a potential map has been knowledge-driven weighted overlay, very
often organized through the analytic hierarchy process. In AHP, the analyst
compares factors pairwise, derives a set of weights from those comparisons,
checks the weights for logical consistency, and then combines the weighted layers
into a single surface. The method has been applied across a remarkable range of
climates and terrains, which is part of why it remains a sensible baseline.

In semi-arid India, Pande et al. \cite{pande2021ahpmif} compared AHP against the
multi-influencing-factor (MIF) technique in the Mula River basin and found AHP
the stronger of the two, with a validation AUC of 0.86 against 0.80 for MIF. Das
and Pardeshi \cite{das2018pravara} took a related but distinct route in the
Pravara basin, using frequency-ratio and influencing-factor statistics over eight
hydrogeological layers and reporting AUC values of 0.73 and 0.69 respectively.
The transferability of these ideas to water-scarce settings closely resembling
the Pothohar is well illustrated by Mallick et al. \cite{mallick2019fuzzyahp},
who applied a fuzzy variant of AHP in the arid Aseer region of Saudi Arabia using
geology, lineaments, geomorphology, NDVI, TWI, slope, and soil. In hard-rock
crystalline terrain---directly relevant to the Salt Range and Kala Chitta
exposures of the study area---Benjmel et al. \cite{benjmel2020morocco} showed
that lineament density and geological structure become the dominant controls,
with fracture zones acting as conduits through otherwise near-impermeable rock.
The same multi-criteria machinery extends naturally to related hazards; Swain et
al. \cite{swain2020floodahp}, for instance, used a cloud-based GIS-AHP workflow
for flood-susceptibility mapping, underlining how general the weighted-overlay
paradigm has become.

These studies establish weighted overlay as a credible, knowledge-grounded
method---which is precisely why it is used here to generate training labels
rather than being discarded. Its limitations are equally well documented,
however. The weights are fixed by expert judgement and are therefore inherently
subjective; the combination rule is linear and additive, which cannot represent
the conditional dependencies that real hydrogeology exhibits; and the output is a
single deterministic surface with no built-in measure of its own accuracy.
Arabameri et al. \cite{arabameri2018shahroud}, comparing statistical,
data-mining, and multi-criteria decision methods on the same plain in Iran, found
the data-driven approaches generally matched or outperformed the knowledge-driven
ones, which points directly toward the machine-learning methods considered next.

\section{From Knowledge-Driven to Machine-Learning Methods}
\label{sec:lit_ml_spatial}

The migration from expert-weighted overlay to learned models has been most
visible in spatial susceptibility mapping---a family of problems, including
landslide, flood, and groundwater mapping, that share the same structure:
predict an ordinal spatial class from a stack of geo-environmental predictors.
Because the structure is shared, methods developed for one transfer readily to
the others, and the landslide-susceptibility literature in particular has become
a proving ground for the tree-based ensembles now common in GWPM.

The most directly analogous study to the present work is that of Sahin
\cite{sahin2020ensembletree}, who compared XGBoost, gradient boosting machine,
and random forest for landslide susceptibility in Turkey using fifteen causative
factors. XGBoost came out ahead, with an overall accuracy of 85.0\% and an AUC
of 0.898, followed by random forest (AUC 0.886) and GBM (AUC 0.880); a Wilcoxon
signed-rank test confirmed the differences were statistically significant. That
ordering---XGBoost ahead of random forest, with gradient boosting close
behind---anticipates the result obtained in this thesis, and lends external
weight to the model ranking reported in Chapter~\ref{ch:results}. Wu et al.
\cite{wu2019adaboost} explored a complementary corner of the ensemble space,
pairing alternating decision trees with AdaBoost and bagging, which provides part
of the rationale for including a voting ensemble among the candidate models here.

Two further studies inform specific design choices. Chang et al.
\cite{chang2019scaleeffects} examined how the resolution of DEM-derived
topographic variables affects model accuracy, comparing five-metre LiDAR,
resampled thirty-metre LiDAR, and thirty-metre ASTER data; the thirty-metre
derivatives performed best, which supports the use of SRTM thirty-metre data for
elevation, slope, aspect, and TWI in this work. They also found non-parametric
models such as random forest and support vector machines to outperform
parametric logistic regression, consistent with the wider trend toward tree-based
methods. On the question of how to split data for honest evaluation, Nguyen et
al. \cite{nguyen2021datasplitting} showed that the train--validation--test
partition materially affects reported performance, which motivates the explicit
72/8/20 split adopted in Chapter~\ref{ch:methodology}.

Perhaps the most important methodological precedent concerns the labels
themselves. Because field measurements are rarely dense enough to cover an entire
study area, several authors have trained models on knowledge-driven labels rather
than on sparse observations. Roy and Saha \cite{roy2019knowledgedriven}
demonstrated that expert-weighted factor combinations---built from rainfall,
slope, aspect, geology, soil, NDVI, and TWI---can serve as reliable spatial
targets, achieving ROC-AUC values around 0.91. This is the same logic that
underpins the weighted-overlay labelling used here, and it is worth being clear
about what it does and does not buy: it provides spatially complete, theoretically
grounded targets, but it also means a model's internal accuracy measures fidelity
to the overlay rather than to ground truth---a point taken up squarely in
Chapters~\ref{ch:results} and~\ref{ch:discussion}.

\section{Tree-Based Learning Algorithms}
\label{sec:lit_algorithms}

The models used in this thesis all belong to the tree-ensemble family, and a
brief note on their lineage is warranted. Random forests build many decision
trees on bootstrapped samples and randomized feature subsets and average their
votes, which reduces variance and resists overfitting; their suitability for
water-resources problems specifically has been reviewed by Tyralis et al.
\cite{tyralis2019randomforests}, building on the widely used implementation
described by Liaw and Wiener \cite{liaw2007randomforest}. Gradient boosting takes
a different route, adding trees sequentially so that each new tree corrects the
errors of those before it. XGBoost \cite{chen2016xgboost} made this approach both
scalable and well regularized through a penalized objective and an efficient,
parallelized tree-construction algorithm, and it has become one of the most
widely used models in applied machine learning. LightGBM \cite{ke2017lightgbm}
pushes efficiency further with histogram-based splitting and leaf-wise tree
growth, which makes it well suited to the large pixel counts encountered in
raster-scale prediction. The choice to compare random forest, XGBoost, LightGBM,
and a soft-voting ensemble therefore spans the two dominant ensemble
philosophies---bagging and boosting---plus their combination.

A recurring theme in this literature is the importance of evaluating such models
with more than a single accuracy figure, particularly when classes are balanced
by construction. Chicco and Jurman \cite{chicco2020mcc} argue for agreement-based
measures over raw accuracy, and the present work accordingly reports macro
$F_1$, Cohen's $\kappa$, and ROC-AUC alongside accuracy.

\section{Machine Learning in Hydrology and Water Resources}
\label{sec:lit_ml_hydrology}

Stepping back from spatial mapping, the past few years have seen machine learning
become a standard tool across hydrology more generally, which provides the wider
context for its use here. Sit et al. \cite{sit2020deeplearning} survey the rapid
adoption of deep-learning methods across the water sciences, from rainfall-runoff
modelling to flood forecasting. Closer to the present topic, Hai et al.
\cite{hai2022gwlreview} review machine-learning models for groundwater-level
prediction specifically, documenting both the breadth of methods in use and the
recurring difficulty of obtaining adequate ground-truth data---the same
constraint that shapes the validation strategy of this thesis. Zhu et al.
\cite{zhu2022mlwaterquality} extend the picture to water-quality evaluation. Read
together, these reviews establish that data-driven modelling is now mainstream in
the field, while also being candid about its dependence on data availability,
which is rarely ideal in the regions where the need is greatest.

\section{Groundwater and Water Resources in the Pothohar Region}
\label{sec:lit_regional}

Bringing the focus back to the study area, a cluster of recent work has
characterized the hydrology and environmental stresses of the \studyarea{} and
the wider Punjab, even if none of it produces an ML groundwater potential map of
the kind developed here.

The closest methodological neighbour is Khan et al. \cite{khan2023rainfallrunoff},
who applied machine-learning models---including stand-alone and wavelet-coupled
neural networks and adaptive neuro-fuzzy systems---to simulate the
rainfall--runoff process in four Pothohar basins, reaching Nash--Sutcliffe
efficiencies as high as 0.92. Their work demonstrates both that data-driven
models perform well in this rainfall-limited setting and that rainfall is the
dominant driver of the region's hydrological response, a finding that resonates
with the feature-importance results reported later in this thesis. The most
spatially comparable study is Hafeez et al. \cite{hafeez2025rwh}, who used AHP to
map rainwater-harvesting suitability across essentially the same five-district
region, identifying roughly 45\% of the area as highly suitable for water storage.
Because recharge-favourable terrain and storage-suitable terrain are governed by
overlapping factors, their suitability map offers a useful independent reference
against which the spatial pattern of this thesis can be sanity-checked, a
comparison developed in Chapter~\ref{ch:discussion}.

The remaining regional literature fills in the surrounding context. Sarwar et al.
\cite{sarwar2022soandrought} document how meteorological droughts in the Soan
basin---the plateau's principal drainage unit---propagate rapidly into
hydrological droughts, underscoring the dependence of the region on groundwater
buffering. Waseem et al. \cite{waseem2022drought} place this in a longer-term
Punjab-wide frame, reporting declining agricultural water availability across
several decades, while a separate review by Waseem et al.
\cite{waseem2024waterquality} highlights shallow-aquifer contamination risk in the
Rawalpindi--Islamabad twin-city area, a reminder that potential and potability are
not the same thing. At the district scale, Khan et al. \cite{khan2024afforestation}
describe the semi-arid, barren-land conditions of Attock, and Ejaz et al.
\cite{ejaz2023soilwater} examine constraints on soil-and-water conservation
practice across the Pothohar. Broader Punjab and Indus-basin studies---land-cover
modelling with random forests \cite{asif2023lulc}, irrigation supply-and-demand
projections \cite{ahmed2021indusbasin}, and urban water security
\cite{kamran2024watersecurity}---complete the regional backdrop and reinforce the
motivation set out in Chapter~\ref{ch:introduction}.

\section{Synthesis and Research Gap}
\label{sec:lit_synthesis}

A few conclusions emerge from this survey. First, the set of geo-environmental
controls on groundwater is well established, and the nine features used here sit
comfortably within it. Second, knowledge-driven weighted overlay and AHP are
proven, defensible methods---good enough to generate training labels---but their
fixed weights, linear combination, and lack of an accuracy estimate are genuine
limitations that data-driven models are designed to relax. Third, tree-based
ensembles, and XGBoost in particular, have repeatedly outperformed both
knowledge-driven and simpler statistical methods on structurally identical
spatial-prediction problems, and the practice of training such models on
knowledge-based labels is itself established. Fourth, machine learning is now firmly
embedded in hydrology, including in the Pothohar region, yet its application there
has been to rainfall-runoff and related problems rather than to groundwater
potential mapping.

The gap is therefore specific and, in light of the literature, somewhat
conspicuous. The methodological shift from manual overlay to learned models has
been completed for landslide and flood susceptibility, and for rainfall-runoff in
the Pothohar itself, but not for groundwater potential across the full plateau.
The reference study by Waqas et al. \cite{waqas2022chakwal} remains confined to a
single district and a manual method with a limited feature set. This thesis sets
out to close that gap: to carry the field's established methodological progression
through to the \studyarea{}, on a richer nine-feature dataset, and---mindful of
the label-circularity caveat raised repeatedly above---to validate the result
against independent hydrogeological evidence rather than internal accuracy alone.
